% Generated by scripts/assemble.py; corrections belong in source/corrections.json.
\examhead{160}{华南师范大学 2025 年数学分析试卷}{152}
% p153-b23; source PDF page 153
\par\Needspace{5\baselineskip}
\blockanchor{p153-b23}
\majorquestion{一}求极限. 每题 10 分, 共 50 分.\par
% p153-b24; source PDF page 153
\blockanchor{p153-b24}
\question{1}$\lim_{n\to \infty}\frac{\sqrt{n + 5} - \sqrt{n}}{\sqrt{n} - \sqrt{n + 2005}}.$\par
% p153-b25; source PDF page 153
\blockanchor{p153-b25}
\question{2}$\lim_{n\to \infty}\left(\frac{n + 2}{n - 1}\right)^n.$\par
% p154-b00; source PDF page 154
\blockanchor{p154-b00}
\question{3}$\lim_{x\to +\infty}\left(\sin \sqrt{x + 1} -\sin \sqrt{x}\right).$\par
% p154-b01; source PDF page 154
\blockanchor{p154-b01}
\question{4}$\lim_{x\to 0}\frac{x\cos x - \sin x}{x^2\sin x}.$\par
% p154-b02; source PDF page 154
\blockanchor{p154-b02}
\question{5}$\lim_{(x,y)\to(0,0)}\frac{xy}{\sqrt{xy+1}-1}.$\par
% p154-b03; source PDF page 154
\par\Needspace{5\baselineskip}
\blockanchor{p154-b03}
\majorquestion{二}计算题. 每题 10 分, 共 40 分.\par
% p154-b04; source PDF page 154
\blockanchor{p154-b04}
\question{6}设 $n$ 为正整数, $y = \frac{x^n}{2 - x}$, 求 $y^{(n)}$.\par
% p154-b05; source PDF page 154
\blockanchor{p154-b05}
\question{7}设 $u = f\left(\frac{y}{x} + \frac{z}{y} + \frac{x}{z}\right)$, 其中 $f$ 二阶可导, 求 $\frac{\partial u}{\partial x}, \frac{\partial^2 u}{\partial x \partial y}$.\par
% p154-b06; source PDF page 154
\blockanchor{p154-b06}
\question{8}求 $\int e^{-x}\left(\frac{1 + x}{1 + x^2}\right)^2\mathrm{d}x.$\par
% p154-b07; source PDF page 154
\blockanchor{p154-b07}
\question{9}计算 $\iint_{D} |\sin(x + y)| \, \mathrm{d}x \, \mathrm{d}y$, $D = [0, \pi] \times [0, \pi]$.\par
% p154-b08; source PDF page 154
\par\Needspace{5\baselineskip}
\blockanchor{p154-b08}
\majorquestion{三}讨论题. 10 分.\par
% p154-b09; source PDF page 154
\par\Needspace{6\baselineskip}
\blockanchor{p154-b09}
\question{10}对于二元函数 $f(x, y)$, 如果两个累次极限存在, 且\par
\nopagebreak[4]
% p154-b10; source PDF page 154
\blockanchor{p154-b10}
\displayformula{\lim _ {x \rightarrow 0} \lim _ {y \rightarrow 0} f (x, y) = \lim _ {y \rightarrow 0} \lim _ {x \rightarrow 0} f (x, y).}
% p154-b11; source PDF page 154
\blockanchor{p154-b11}
问 $\lim_{(x,y)\to (0,0)}f(x,y)$ 是否一定存在？请说明理由.\par
% p154-b12; source PDF page 154
\par\Needspace{5\baselineskip}
\blockanchor{p154-b12}
\majorquestion{四}解答题.\par
% p154-b13; source PDF page 154
\blockanchor{p154-b13}
\question{11}(10 分) 设 $f_{n}(x) = \frac{x^{2}}{1 + n^{3}x^{3}}$, $n = 1,2,\dots$, 试证明函数列 $\{f_{n}(x)\}$ 在 $(0, +\infty)$ 内一致收敛.\par
% p154-b14; source PDF page 154
\blockanchor{p154-b14}
\question{12}(20 分) 判断下列级数或反常积分的敛散性.\par
% p154-b15; source PDF page 154
\blockanchor{p154-b15}
(1) $\sum_{n=1}^{\infty} \frac{1}{n} \sin \frac{1}{n^q}, q \geq 0.$\par
% p154-b16; source PDF page 154
\blockanchor{p154-b16}
(2) $\int_{3}^{+\infty}\frac{\ln x\mathrm{d}x}{(x - \sqrt{x} + 2)^q},q > 0.$\par
% p154-b17; source PDF page 154
\blockanchor{p154-b17}
\question{13}(10 分) 求函数 $f(x) = \frac{4}{x^2 - 2x - 3}$ 在 $x = 2$ 处的幂级数展开式.\par
% p154-b18; source PDF page 154
\blockanchor{p154-b18}
\question{14}(10 分) 设在矩形区域 $D = [a, b] \times [c, d]$ 上, $F(x, y)$ 关于 $x$ 连续, $F_y(x, y)$ 存在且有界. 证明: $F(x, y)$ 在 $D$ 上连续.\par
